\documentclass[twocolumn]{article}

\usepackage{gvv}
\usepackage{amsmath}
\usepackage{amsfonts}
\usepackage{tabularx}
\usepackage{listings} % Required for lstlisting to work
\usepackage{graphicx} % <-- make sure this is in the preamble!

\begin{document}

% Your title
\title{GATE EE 2025 Questions}
\author{}
\date{}
\maketitle

% Now start questions


% --------------------------
% QUESTION 1
% --------------------------
1. Identify the CORRECT statement for the following vectors $\vec{a} = 3\hat{i} + 2\hat{j}$ and $\vec{b} = 2\hat{i} + \hat{j}$.


\begin{itemize}
\item[(A)] The vectors $\vec{a}$ and $\vec{b}$ are linearly independent.

\item[(B)] The vectors $\vec{a}$ and $\vec{b}$ are linearly dependent.

\item[(C)] The vectors $\vec{a}$ and $\vec{b}$ are orthogonal.

\item[(D)] The vectors $\vec{a}$ and $\vec{b}$ are normalized.

(GATE EE 2025)
  \end{itemize}




% --------------------------
% QUESTION 2
% --------------------------
2. Two uniform thin rods of equal length $L$ and masses $M_1$ and $M_2$ are joined together along the length. The moment of inertia of the combined rod of length $2L$ about an axis passing through the mid-point and perpendicular to the length of the rod is,

\begin{itemize}
\item[(A)] $\frac{1}{12} L^2 (M_1 + M_2)$

\item[(B)] $\frac{1}{6} L^2 (M_1 + M_2)$

\item[(C)] $\frac{1}{3} L^2 (M_1 + M_2)$

\item[(D)] $\frac{1}{2} L^2 (M_1 + M_2)$

(GATE EE 2025)
  \end{itemize}




% --------------------------
% QUESTION 3
% --------------------------
3. The space-time dependence of the electric field of a linearly polarized light in free space is given by $E_x = E_0 \cos(\omega t - kz)$ where $E_0, \omega$ and $k$ are the amplitude, the angular frequency and the wavevector respectively. The time-averaged energy density associated with the electric field is

\begin{itemize}
\item[(A)] $\frac{1}{4} \epsilon_0 E_0^2$

\item[(B)] $\frac{1}{2} \epsilon_0 E_0^2$

\item[(C)] $\epsilon_0 E_0^2$

\item[(D)] $2 \epsilon_0 E_0^2$

(GATE EE 2025)
  \end{itemize}




% --------------------------
% QUESTION 4
% --------------------------
4. If the peak output voltage of a full wave rectifier is 10 V, its d.c. voltage is

\begin{itemize}
\item[(A)] 10.0 V

\item[(B)] 7.07 V

\item[(C)] 6.36 V

\item[(D)] 3.18 V

(GATE EE 2025)
  \end{itemize}




% --------------------------
% QUESTION 5
% --------------------------
5. A particle of mass $m$ is confined in a two-dimensional square well potential of dimension $a$.  
This potential $V(x,y)$ is given by:


\[
V(x,y) = 0 \quad \text{for} \quad -a < x < a,\; -a < y < a
\quad \text{and} \quad
V(x,y) = \infty \; \text{elsewhere}.
\]

The energy of the first excited state for this particle is given by

\begin{itemize}
\item[(A)] $\frac{\pi^2 \hbar^2}{2ma^2}$

\item[(B)] $\frac{2 \pi^2 \hbar^2}{2ma^2}$

\item[(C)] $\frac{5 \pi^2 \hbar^2}{2ma^2}$

\item[(D)] $\frac{4 \pi^2 \hbar^2}{2ma^2}$

(GATE EE 2025)
  \end{itemize}




% --------------------------
% QUESTION 6
% --------------------------
6. The isothermal compressibility, $\kappa$ of an ideal gas at temperature $T_0$ and volume $V_0$ is given by

\begin{itemize}
\item[(A)] $\displaystyle -\frac{1}{V_0} \left( \frac{\partial V}{\partial P} \right)_{T_0}$

\item[(B)] $\displaystyle \frac{1}{V_0} \left( \frac{\partial V}{\partial P} \right)_{T_0}$

\item[(C)] $\displaystyle -\frac{1}{V_0} \left( \frac{\partial P}{\partial V} \right)_{T_0}$

\item[(D)] $\displaystyle \frac{1}{V_0} \left( \frac{\partial P}{\partial V} \right)_{T_0}$

(GATE EE 2025)
  \end{itemize}




% --------------------------
% QUESTION 7
% --------------------------
7. The ground state of sodium atom ($^{11}$Na) is a ${}^2S_{1/2}$ state.  
The difference in energy levels arising in the presence of a weak external magnetic field $B$, given in terms of Bohr magneton $\mu_B$, is

\begin{itemize}
\item[(A)] $\mu_B B$ \\
\item[(B)] $2 \mu_B B$ \\
\item[(C)] $4 \mu_B B$ \\
\item[(D)] $6 \mu_B B$

(GATE EE 2025)
  \end{itemize}




% --------------------------
% QUESTION 8
% --------------------------
8. For an ideal Fermi gas in three dimensions, the electron velocity $v_F$ at the Fermi surface is related to electron concentration $n$ as

\begin{itemize}
\item[(A)] $v_F \propto n^{2/3}$

\item[(B)] $v_F \propto n$

\item[(C)] $v_F \propto n^{1/2}$

\item[(D)] $v_F \propto n^{1/3}$

(GATE EE 2025)
  \end{itemize}




% --------------------------
% QUESTION 9
% --------------------------
9. Which one of the following sets corresponds to fundamental particles?

\begin{itemize}
\item[(A)] proton, electron and neutron

\item[(B)] proton, electron and photon

\item[(C)] electron, photon and neutrino

\item[(D)] quark, electron and meson

(GATE EE 2025)
  \end{itemize}




% --------------------------
% QUESTION 10
% --------------------------
10. In case of a Geiger-Muller (GM) counter, which one of the following statements is CORRECT?

\begin{itemize}
\item[(A)] Multiplication factor of the detector is of the order of $10^{10}$

\item[(B)] Type of the particles detected can be identified

\item[(C)] Energy of the particles detected can be distinguished

\item[(D)] Operating voltage of the detector is few tens of Volts

(GATE EE 2025)
  \end{itemize}




% --------------------------
% QUESTION 11
% --------------------------
11. A plane electromagnetic wave traveling in free space is incident normally on a glass plate of refractive index $n = \frac{3}{2}$.  
If there is no absorption by the glass, its reflectivity is

\begin{itemize}
\item[(A)] 4\% \\
\item[(B)] 16\% \\
\item[(C)] 20\% \\
\item[(D)] 50\%

(GATE EE 2025)
  \end{itemize}




% --------------------------
% QUESTION 12
% --------------------------
12. A Ge semiconductor is doped with acceptor impurity concentration of $10^{15}$ atoms/cm$^3$.  
For the given hole mobility of 1800 cm$^2$/V-s, the resistivity of this material is

\begin{itemize}
\item[(A)] 0.288 $\Omega$ cm

\item[(B)] 0.694 $\Omega$ cm

\item[(C)] 3.472 $\Omega$ cm

\item[(D)] 6.944 $\Omega$ cm

(GATE EE 2025)
  \end{itemize}




% --------------------------
% QUESTION 13
% --------------------------
13. A classical gas of molecules, each of mass $m$, is in thermal equilibrium at the absolute temperature $T$.  
The velocity components of the molecules along the Cartesian axes are $v_x$, $v_y$, and $v_z$.  
The mean value of $(v_x^2 + v_y^2)$ is

\begin{itemize}
\item[(A)] $\frac{k_B T}{m}$

\item[(B)] $\frac{3}{2} \frac{k_B T}{m}$

\item[(C)] $\frac{1}{2} \frac{k_B T}{m}$

\item[(D)] $2 \frac{k_B T}{m}$

(GATE EE 2025)
  \end{itemize}




% --------------------------
% QUESTION 14
% --------------------------
14. In a central force field, the trajectory of a particle of mass $m$ and angular momentum $L$ in plane polar coordinates is given by


\[
\frac{1}{r} = \frac{m}{L} \big( 1 + \varepsilon \cos \theta \big)
\]

where $\varepsilon$ is the eccentricity of the particle’s motion. Which one of the following choices for $\varepsilon$ gives rise to a parabolic trajectory?

\begin{itemize}
\item[(A)] $\varepsilon = 0$ \\
\item[(B)] $\varepsilon = 1$ \\
\item[(C)] $0 < \varepsilon < 1$ \\
\item[(D)] $\varepsilon > 1$

(GATE EE 2025)
  \end{itemize}






% --------------------------
% QUESTION 15
% --------------------------
15. Identify the CORRECT energy band diagram for Silicon doped with Arsenic.  
Here CB, VB, $E_D$ and $E_F$ are the conduction band, valence band, impurity level and Fermi level, respectively.




\includegraphics[width=0.9\linewidth]{B_15.png}

(GATE EE 2025)





% --------------------------
% QUESTION 16
% --------------------------
16. The first Stokes line of a rotational Raman spectrum is observed at 12.96 cm$^{-1}$.  
Considering the rigid rotor approximation, the rotational constant is given by

\begin{itemize}
\item[(A)] 6.48 cm$^{-1}$

\item[(B)] 3.24 cm$^{-1}$

\item[(C)] 2.16 cm$^{-1}$

\item[(D)] 1.62 cm$^{-1}$

(GATE EE 2025)
  \end{itemize}




% --------------------------
% QUESTION 17
% --------------------------
17. The total energy, $E$ of an ideal non-relativistic Fermi gas in three dimensions is given by

\[
E \propto \left( \frac{N}{V} \right)^{5/3}
\]

where $N$ is the number of particles and $V$ is the volume of the gas.  
Identify the CORRECT equation of state (with $P$ being the pressure):

\begin{itemize}
\item[(A)] $PV = \frac{1}{3} E$

\item[(B)] $PV = \frac{2}{3} E$

\item[(C)] $PV = E$

\item[(D)] $PV = \frac{5}{3} E$

(GATE EE 2025)
  \end{itemize}




% --------------------------
% QUESTION 18
% --------------------------
18. Consider the wavefunction

\[
\Psi = \psi(\mathbf{r}_1, \mathbf{r}_2) \; \chi_s
\]

for a fermionic system consisting of two spin-half particles.  
The spatial part of the wavefunction is given by

\[
\psi(\mathbf{r}_1, \mathbf{r}_2) = \frac{1}{2} 
\Big[ \phi_1(\mathbf{r}_1) \phi_2(\mathbf{r}_2) + \phi_1(\mathbf{r}_2) \phi_2(\mathbf{r}_1) \Big].
\]

The spin part $\chi_s$ of the wavefunction with spin states $\alpha$ ($+\tfrac{1}{2}$) and $\beta$ ($-\tfrac{1}{2}$) should be

\begin{itemize}
\item[(A)] $\frac{1}{\sqrt{2}} (\alpha \beta + \beta \alpha)$

\item[(B)] $\frac{1}{\sqrt{2}} (\alpha \beta - \beta \alpha)$

\item[(C)] $\alpha \alpha$

\item[(D)] $\beta \beta$

(GATE EE 2025)
  \end{itemize}




% --------------------------
% QUESTION 19
% --------------------------
19. The electric and the magnetic fields $E(z,t)$ and $B(z,t)$ respectively, corresponding to the scalar potential $\phi(z,t) = 0$ and the vector potential $A(z,t) = itz$ are

\begin{itemize}
\item[(A)] $E = i z$ and $B = -j t$

\item[(B)] $E = i z$ and $B = j t$

\item[(C)] $E = -i z$ and $B = -j t$

\item[(D)] $E = -i z$ and $B = j t$

(GATE EE 2025)
  \end{itemize}




% --------------------------
% QUESTION 20
% --------------------------
20. Consider the following OP-AMP circuit.


\includegraphics[width=0.6\linewidth]{A_20.png}


Which one of the following correctly represents the output $V_{out}$ corresponding to the input $V_{in}$?


\includegraphics[width=0.6\linewidth]{A_20.1.png}
(GATE EE 2025)





% --------------------------
% QUESTION 21
% --------------------------
21. Deuteron has only one bound state with spin parity $1^+$, isospin 0 and electric quadrupole moment 0.286 efm$^2$.  
These data suggest that the nuclear forces are having


\begin{itemize}
\item[(A)] only spin and isospin dependence

\item[(B)] no spin dependence and no tensor components

\item[(C)] spin dependence but no tensor components

\item[(D)] spin dependence along with tensor components

(GATE EE 2025)
  \end{itemize}




% --------------------------
% QUESTION 22
% --------------------------
22. A particle of unit mass moves along the $x$-axis under the influence of a potential,

\[
V(x) = x^2 - 2x.
\]

The particle is found to be in stable equilibrium at the point $x = 2$.  
The time period of oscillation of the particle is

\begin{itemize}
\item[(A)] $\frac{2}{\pi}$

\item[(B)] $\pi$

\item[(C)] $\frac{3}{2} \pi$

\item[(D)] $2\pi$

(GATE EE 2025)
  \end{itemize}




% --------------------------
% QUESTION 23
% --------------------------
23. Which one of the following CANNOT be explained by considering a harmonic approximation for the lattice vibrations in solids?

\begin{itemize}
\item[(A)] Debye’s $T^3$ law

\item[(B)] Dulong–Petit’s law

\item[(C)] Optical branches in lattices

\item[(D)] Thermal expansion

(GATE EE 2025)
  \end{itemize}




% --------------------------
% QUESTION 24
% --------------------------
24. A particle is constrained to move in a truncated harmonic potential well ($x > 0$) as shown in the figure.  
Which one of the following statements is CORRECT?


\includegraphics[width=0.6\linewidth]{A_24.png}

\begin{itemize}
\item[(A)] The parity of the first excited state is even

\item[(B)] The parity of the ground state is even

\item[(C)] The ground state energy is $\frac{1}{2} \omega$

\item[(D)] The first excited state energy is $\frac{7}{2} \omega$

(GATE EE 2025)
  \end{itemize}




% --------------------------
% QUESTION 25
% --------------------------
25. The number of independent components of the symmetric tensor $A_{ij}$ with indices $i,j = 1,2,3$ is

\begin{itemize}
\item[(A)] 1

\item[(B)] 3

\item[(C)] 6

\item[(D)] 9

(GATE EE 2025)
  \end{itemize}




% --------------------------
% QUESTION 26
% --------------------------
26. Consider a system in the unperturbed state described by the Hamiltonian,



\[
H_0 = \begin{pmatrix}
1 & 0 \\
0 & 1
\end{pmatrix}.
\]

The system is subjected to a perturbation of the form


\[
H' = \begin{pmatrix}
\delta & \delta \\
\delta & \delta
\end{pmatrix},
\quad \text{where } \delta \ll 1.
\]

The energy eigenvalues of the perturbed system using the first order perturbation approximation are
\begin{itemize}
\item[(A)] $1$ and $(1 + 2\delta)$

\item[(B)] $(1 + \delta)$ and $(1 - \delta)$

\item[(C)] $(1 + 2\delta)$ and $(1 - 2\delta)$

\item[(D)] $(1 + \delta)$ and $(1 - 2\delta)$

(GATE EE 2025)
  \end{itemize}




% --------------------------
% QUESTION 27
% --------------------------
27. Inverse susceptibility ($1/\chi$) as a function of temperature $T$ for a material undergoing a paramagnetic to ferromagnetic transition is given in the figure, where $O$ is the origin.  
The values of the Curie constant $C$ and the Weiss molecular field constant $\lambda$, in CGS units, are


\includegraphics[width=0.6\linewidth]{A_27.png}

\begin{itemize}
\item[(A)] $C = 5 \times 10^{-5}, \; \lambda = 3 \times 10^{-2}$

\item[(B)] $C = 3 \times 10^{-2}, \; \lambda = 5 \times 10^{-5}$

\item[(C)] $C = 3 \times 10^{-2}, \; \lambda = 2 \times 10^{-4}$

\item[(D)] $C = 2 \times 10^{-4}, \; \lambda = 3 \times 10^{-2}$

(GATE EE 2025)
  \end{itemize}




% --------------------------
% QUESTION 28
% --------------------------
28. A plane polarized electromagnetic wave in free space at time $t = 0$ is given by


\[
E(x,z) = 10\, \hat{j} + \hat{i} \; \exp[ i(6x + 8z) ].
\]

The magnetic field $B(x,z,t)$ is given by

\begin{itemize}
\item[(A)] $\displaystyle B(x,z,t) = \frac{1}{c} \; \hat{k} \times (\hat{i} + \hat{j}) \; \exp[ i(6x + 8z - 10ct) ]$

\item[(B)] $\displaystyle B(x,z,t) = \frac{1}{c} \; \hat{k} \times (\hat{i} - \hat{j}) \; \exp[ i(6x + 8z - 10ct) ]$

\item[(C)] $\displaystyle B(x,z,t) = \frac{1}{c} \; \hat{k} \times (\hat{i} + \hat{j}) \; \exp[ i(6x + 8z) ]$

\item[(D)] $\displaystyle B(x,z,t) = \frac{1}{c} \; \hat{k} \times (\hat{i} - \hat{j}) \; \exp[ i(6x + 8z) ]$

(GATE EE 2025)
  \end{itemize}




% --------------------------
% QUESTION 29
% --------------------------
29. The eigenvalues of the matrix


\[
\begin{pmatrix}
0 & 1 & 0 \\
1 & 0 & 1 \\
0 & 1 & 0
\end{pmatrix}
\]

are

\begin{itemize}
\item[(A)] $0, 1, 1$

\item[(B)] $0, 2, -2$

\item[(C)] $\frac{1}{2}, -\frac{1}{2}, 0$

\item[(D)] $2, -2, 0$

(GATE EE 2025)
  \end{itemize}




% --------------------------
% QUESTION 30
% --------------------------
30. Match the typical spectroscopic regions specified in Group I with the corresponding type of transitions in Group II.

\begin{center}
\begin{tabularx}{\linewidth}{lX}
Group I & Group II \\

(P) Infra-red region & (i) electronic transitions involving valence electrons \\
(Q) Ultraviolet-visible region & (ii) nuclear transitions \\
(R) X-ray region & (iii) vibrational transitions of molecules \\
(S) $\gamma$-ray region & (iv) transitions involving inner shell electrons \\

\end{tabularx}
\end{center}

\begin{itemize}
\item[(A)] (P, i); (Q, iii); (R, ii); (S, iv)
\item[(B)] (P, ii); (Q, iv); (R, i); (S, iii)
\item[(C)] (P, iii); (Q, i); (R, iv); (S, ii)
\item[(D)] (P, iv); (Q, i); (R, ii); (S, iii)
\end{itemize}

(GATE EE 2025)





% --------------------------
% QUESTION 31
% --------------------------
31. In the following circuit, for the output voltage to be $V_o = \left( -V_1 + \frac{V_2}{2} \right)$, the ratio $R_1 / R_2$ is


\includegraphics[width=0.6\linewidth]{A_31.png}

\begin{itemize}
\item[(A)] $\frac{1}{2}$

\item[(B)] 1

\item[(C)] 2

\item[(D)] 3

(GATE EE 2025)
  \end{itemize}




% --------------------------
% QUESTION 32
% --------------------------
32. The terms $\{ j_1, j_2 \}_J$ arising from $2s^1 3d^1$ electronic configuration in the $j$–$j$ coupling scheme are

\begin{itemize}
\item[(A)] $\left\{ \frac{1}{2}, 1 \right\}_{\frac{3}{2}}, \; \left\{ \frac{1}{2}, 2 \right\}_{\frac{3}{2}}, \; \left\{ \frac{1}{2}, 2 \right\}_{\frac{5}{2}}$

\item[(B)] $\left\{ \frac{1}{2}, 0 \right\}_{\frac{1}{2}}, \; \left\{ \frac{1}{2}, 1 \right\}_{\frac{3}{2}}$

\item[(C)] $\left\{ \frac{1}{2}, 0 \right\}_{\frac{1}{2}}, \; \left\{ \frac{1}{2}, 2 \right\}_{\frac{3}{2}}, \; \left\{ \frac{1}{2}, 2 \right\}_{\frac{5}{2}}$

\item[(D)] $\left\{ \frac{1}{2}, 1 \right\}_{\frac{3}{2}}, \; \left\{ \frac{1}{2}, 2 \right\}_{\frac{3}{2}}, \; \left\{ \frac{1}{2}, 2 \right\}_{\frac{5}{2}}$

(GATE EE 2025)
  \end{itemize}




% --------------------------
% QUESTION 33
% --------------------------
33. In the following circuit, the voltage drop across the ideal diode in forward bias condition is 0.7 V.  
The current passing through the diode is


\includegraphics[width=0.6\linewidth]{A_31.png}
\begin{itemize}
\item[(A)] 0.5 mA

\item[(B)] 1.0 mA

\item[(C)] 1.5 mA

\item[(D)] 2.0 mA

(GATE EE 2025)
  \end{itemize}




% --------------------------
% QUESTION 34
% --------------------------
34. Choose the CORRECT statement from the following.
\begin{itemize}
\item[(A)] Neutron interacts through electromagnetic interaction

\item[(B)] Electron does not interact through weak interaction

\item[(C)] Neutrino interacts through weak and electromagnetic interaction

\item[(D)] Quark interacts through strong interaction but not through weak interaction

  \end{itemize}

  (GATE EE 2025)





% --------------------------
% QUESTION 35
% --------------------------
35. A rod of proper length $l_0$ oriented parallel to the $x$-axis moves with speed $\frac{2c}{3}$ along the $x$-axis in the $S$-frame, where $c$ is the speed of light in free space.  
The observer is also moving along the $x$-axis with speed $\frac{c}{2}$ with respect to the $S$-frame.  
The length of the rod as measured by the observer is

\begin{itemize}
\item[(A)] 0.35 $l_0$

\item[(B)] 0.48 $l_0$

\item[(C)] 0.87 $l_0$

\item[(D)] 0.97 $l_0$

(GATE EE 2025)
  \end{itemize}




% --------------------------
% QUESTION 36
% --------------------------
36. A simple cubic crystal with lattice parameter $a_c$ undergoes transition into a tetragonal structure with lattice parameters $a_t = b_t = \sqrt{2} a_c$ and $c_t = 2a_c$, below a certain temperature.  
The ratio of the interplanar spacings of $(101)$ planes for the cubic and the tetragonal structures is

\begin{itemize}
\item[(A)] $\frac{1}{\sqrt{6}}$

\item[(B)] $\frac{1}{6}$

\item[(C)] $\frac{3}{8}$

\item[(D)] $\frac{3}{\sqrt{8}}$

(GATE EE 2025)
  \end{itemize}




% --------------------------
% QUESTION 37
% --------------------------
37. Consider the following circuit in which the current gain $\beta_{dc}$ of the transistor is 100.

% Insert the transistor circuit diagram here
\includegraphics[width=0.6\linewidth]{A_37.png}

Which one of the following correctly represents the load line (collector current $I_C$ with respect to collector-emitter voltage $V_{CE}$) and Q-point of this circuit?

\includegraphics[width=0.6\linewidth]{B_37.png}
(GATE EE 2025)





% --------------------------
% QUESTION 38
% --------------------------
38. Consider a system whose three energy levels are given by $0$, $\varepsilon$ and $2\varepsilon$.  
The energy level $\varepsilon$ is two-fold degenerate and the other two are non-degenerate.  
The partition function of the system with $\beta = 1/(k_B T)$ is given by

\begin{itemize}
\item[(A)] $1 + 2 e^{-\beta \varepsilon}$

\item[(B)] $1 + e^{-\beta \varepsilon} + e^{-2\beta \varepsilon}$

\item[(C)] $2 (1 + e^{-\beta \varepsilon})$

\item[(D)] $1 + e^{-\beta \varepsilon} + 2 e^{-2\beta \varepsilon}$

(GATE EE 2025)
  \end{itemize}




% --------------------------
% QUESTION 39
% --------------------------
39. Two infinitely extended homogeneous isotropic dielectric media (medium-1 and medium-2 with dielectric constants $\epsilon_1 = 2\epsilon_0$ and $\epsilon_2 = 5\epsilon_0$, respectively) meet at the $z = 0$ plane as shown in the figure.  
A uniform electric field exists everywhere. For $z \geq 0$, the electric field is given by

\[
\vec{E}_1 = \hat{i} - 2\hat{j} + 3\hat{k}.
\]

The interface separating the two media is charge free.  
The electric displacement vector in medium-2 is given by


\includegraphics[width=0.6\linewidth]{A_39.png}
(GATE EE 2025)





% --------------------------
% QUESTION 40
% --------------------------
40. The ground state wavefunction for the hydrogen atom is given by

\[
\psi_{100} = \frac{1}{\sqrt{\pi a_0^3}} \, e^{-r/a_0},
\]

where $a_0$ is the Bohr radius.

The plot of the radial probability density $P(r)$ for the hydrogen atom in the ground state is
\includegraphics[width=0.6\linewidth]{A_40.png}

(GATE EE 2025)
 
% --------------------------
% QUESTION 41
% --------------------------
41. Total binding energies of $^{15}$O, $^{16}$O and $^{17}$O are 111.96 MeV, 127.62 MeV and 131.76 MeV, respectively.  
The energy gap between $1p_{1/2}$ and $1d_{5/2}$ neutron shells for nuclei whose mass number is close to 16 is


\begin{itemize}
\item[(A)] 4.1 MeV

\item[(B)] 11.5 MeV

\item[(C)] 15.7 MeV

\item[(D)] 19.8 MeV

(GATE EE 2025)
  \end{itemize}




% --------------------------
% QUESTION 42
% --------------------------
42. A particle of mass $m$ is attached to a fixed point $O$ by a weightless inextensible string of length $a$.  
It is rotating under gravity as shown in the figure.

The Lagrangian of the particle is

\[
L = \frac{1}{2} m a^2 \Big( \dot{\theta}^2 + \dot{\phi}^2 \sin^2 \theta \Big) + mga \cos \theta,
\quad \text{where } \theta \text{ and } \phi \text{ are the polar angles.}
\]

\includegraphics[width=0.6\linewidth]{A_27.png}

The Hamiltonian $H$ is

\begin{itemize}
  \item[(A)] $H = \frac{1}{2ma^2} \left( p_\theta^2 + \frac{p_\phi^2}{\sin^2 \theta} - 2 p_\theta p_\phi \cos \theta \right) + mga \cos \theta$
  \item[(B)] $H = \frac{1}{2ma^2} \left( p_\theta^2 + \frac{p_\phi^2}{\sin^2 \theta} + 2 p_\theta p_\phi \cos \theta \right) + mga \cos \theta$
  \item[(C)] $H = \frac{1}{2ma^2} \Big( p_\theta^2 + p_\phi^2 - 2 p_\theta p_\phi \cos \theta \Big) + mga \cos \theta$
  \item[(D)] $H = \frac{1}{2ma^2} \Big( p_\theta^2 + p_\phi^2 + 2 p_\theta p_\phi \cos \theta \Big) + mga \cos \theta$
\end{itemize}

(GATE EE 2025)





% --------------------------
% QUESTION 43
% --------------------------
43. Given
\[
\vec{F} = \vec{r} \times \vec{B}
\quad \text{and} \quad
\vec{B} = B_0 (\hat{i} + \hat{j} + \hat{k}).
\]

Here, $\vec{B}$ is a constant vector and $\vec{r}$ is the position vector.

The value of the line integral
\[
\oint_C \vec{F} \cdot d\vec{r}
\]
where $C$ is a circle of unit radius centered at the origin, is

 %

\includegraphics[width=0.9\linewidth]{A_43.png}






% --------------------------
% QUESTION 44
% --------------------------
44. The value of the integral

\[
\oint_C \frac{e^z}{z} \, dz
\]

using the contour \( C \) of a circle with unit radius \( |z| = 1 \), is

\begin{itemize}
\item[(A)] 0

\item[(B)] $1 - 2\pi i$

\item[(C)] $1 + 2\pi i$

\item[(D)] $2\pi i$
\end{itemize}

(GATE EE 2025)




% --------------------------
% QUESTION 45
% --------------------------
45. A paramagnetic system consisting of $N$ spin-half particles is placed in an external magnetic field.  
It is found that $N/2$ spins are aligned parallel and the remaining $N/2$ spins are aligned antiparallel to the magnetic field.  
The statistical entropy of the system is

\begin{itemize}
\item[(A)] $2 k_B \ln 2$

\item[(B)] $\frac{N}{2} k_B \ln 2$

\item[(C)] $\frac{3}{2} N k_B \ln 2$

\item[(D)] $ N k_B \ln 2$

(GATE EE 2025)
  \end{itemize}





% --------------------------
% QUESTION 46
% --------------------------
46. The equilibrium vibration frequency for an oscillator is observed at 2990 cm$^{-1}$.  
The ratio of the frequencies corresponding to the first and the fundamental spectral lines is 1.96.  
Considering the oscillator to be anharmonic, the anharmonicity constant is

\begin{itemize}
\item[(A)] 0.005

\item[(B)] 0.02

\item[(C)] 0.05

\item[(D)] 0.1

(GATE EE 2025)
  \end{itemize}




% --------------------------
% QUESTION 47
% --------------------------
47. At a certain temperature $T$, the average speed of nitrogen molecules in air is found to be 400 m/s.  
The most probable and the root mean square speeds of the molecules are, respectively,

\begin{itemize}
\item[(A)] 355 m/s, 434 m/s

\item[(B)] 820 m/s, 917 m/s

\item[(C)] 152 m/s, 301 m/s

\item[(D)] 422 m/s, 600 m/s

(GATE EE 2025)
  \end{itemize}




% --------------------------
% COMMON DATA FOR Q48 AND Q49
% --------------------------
\textbf{Common Data for Questions 48 and 49:}

The wavefunction of a particle moving in free space is given by

\[
\psi = e^{i k x} + e^{-i k x}.
\]



% --------------------------
% QUESTION 48
% --------------------------
48. The energy of the particle is

\begin{itemize}
\item[(A)] $\frac{5 \hbar^2 k^2}{2m}$

\item[(B)] $\frac{4 \hbar^2 k^2}{3m}$

\item[(C)] $\frac{\hbar^2 k^2}{2m}$

\item[(D)] $\frac{\hbar^2 k^2}{m}$

(GATE EE 2025)
  \end{itemize}




% --------------------------
% QUESTION 49
% --------------------------
49. The probability current density for the real part of the wavefunction is

\begin{itemize}
\item[(A)] 1

\item[(B)] $\frac{k}{m}$

\item[(C)] $\frac{2k}{m}$

\item[(D)] 0

(GATE EE 2025)
  \end{itemize}




% --------------------------
% COMMON DATA FOR Q50 AND Q51
% --------------------------
\noindent
\textbf{Common Data for Questions 50 and 51:}

The dispersion relation for a one-dimensional monatomic crystal with lattice spacing $a$, which interacts via nearest neighbour harmonic potential, is given by

\[
\omega = A \sin \left( \frac{Ka}{2} \right)
\]

where $A$ is a constant of appropriate unit.




% --------------------------
% QUESTION 50
% --------------------------
50. The group velocity at the boundary of the first Brillouin zone is

\begin{itemize}
\item[(A)] 0

\item[(B)] 1

\item[(C)] $\frac{\sqrt{2}}{2} Aa$

\item[(D)] $\frac{1}{2\sqrt{2}} Aa$

(GATE EE 2025)
  \end{itemize}




% --------------------------
% QUESTION 51
% --------------------------
51. The force constant between the nearest neighbours of the lattice is ($M$ is the mass of the atom)

\begin{itemize}
\item[(A)] $4 M A^2$

\item[(B)] $2 M A^2$

\item[(C)] $M A^2$

\item[(D)] $2 M A^2$

(GATE EE 2025)
  \end{itemize}




% --------------------------
% LINKED ANSWER Q52–Q53
% --------------------------
\textbf{Linked Answer Questions}

\textbf{Statement for Linked Answer Questions 52 and 53:}

In a hydrogen atom, consider that the electronic charge is uniformly distributed in a spherical volume of radius $a$ ($= 0.5 \times 10^{-10}$ m) around the proton.  
The atom is placed in a uniform electric field $E = 30 \times 10^{5}$ V/m.  
Assume that the spherical distribution of the negative charge remains undistorted under the electric field.



52. In the equilibrium condition, the separation between the positive and the negative charge centers is

\begin{itemize}
\item[(A)] $8.66 \times 10^{-16}$ m

\item[(B)] $2.60 \times 10^{-15}$ m

\item[(C)] $2.60 \times 10^{-16}$ m

\item[(D)] $8.66 \times 10^{-15}$ m

(GATE EE 2025)
  \end{itemize}




53. The polarizability of the hydrogen atom in units of (C$^2$m/N) is

\begin{itemize}
\item[(A)] $2.0 \times 10^{-40}$

\item[(B)] $1.4 \times 10^{-41}$

\item[(C)] $1.4 \times 10^{-40}$

\item[(D)] $2.0 \times 10^{-39}$

(GATE EE 2025)
  \end{itemize}




% --------------------------
% LINKED ANSWER Q54–Q55
% --------------------------
\textbf{Statement for Linked Answer Questions 54 and 55:}

A particle of mass $m$ slides under gravity without friction along the parabolic path $y = a x^2$ as shown in the figure. Here $a$ is a constant.

\includegraphics[width=0.6\linewidth]{A_53.png}



54. The Lagrangian for this particle is given by

\begin{itemize}
\item[(A)] $L = \frac{1}{2} m \dot{x}^2 - mg a x$

\item[(B)] $L = \frac{1}{2} m (1 + 4 a^2 x^2) \dot{x}^2 - mg a x^2$

\item[(C)] $L = \frac{1}{2} m \dot{x}^2 + mg a x$

\item[(D)] $L = \frac{1}{2} m (1 + 4 a^2 x^2) \dot{x}^2 + mg a x^2$

(GATE EE 2025)
  \end{itemize}





55. The Lagrange’s equation of motion of the particle is

\begin{itemize}
\item[(A)] $\ddot{x} = 2 g a x$

\item[(B)] $m (1 + 4 a^2 x^2) \ddot{x} + 2 m a x \dot{x}^2 + mg a x^2 = 0$

\item[(C)] $m (1 + 4 a^2 x^2) \ddot{x} + 2 m a x \dot{x}^2 + mg a x^2 = mg a x^2$

\item[(D)] $\ddot{x} = -2 g a x$

(GATE EE 2025)
  \end{itemize}




% --------------------------
% GENERAL APTITUDE Q56–Q60
% --------------------------
\textbf{General Aptitude (GA) Questions}

56. Choose the grammatically INCORRECT sentence:

\begin{itemize}
\item[(A)] They gave us the money back less the service charges of Three Hundred rupees.

\item[(B)] This country’s expenditure is not less than that of Bangladesh.

\item[(C)] The committee initially asked for a funding of Fifty Lakh rupees, but later settled for a lesser sum.

\item[(D)] This country’s expenditure on educational reforms is very less.

(GATE EE 2025)
  \end{itemize}




57. Which one of the following options is the closest in meaning to the word given below?  
\textbf{Mitigate}

\begin{itemize}
\item[(A)] Diminish

\item[(B)] Divulge

\item[(C)] Dedicate

\item[(D)] Denote

  \end{itemize}



58. Choose the most appropriate alternative from the options given below to complete the following sentence:  
\textit{Despite several \rule{2cm}{0.15mm} the mission succeeded in its attempt to resolve the conflict.}


\begin{itemize}
\item[(A)] attempts

\item[(B)] setbacks

\item[(C)] meetings

\item[(D)] delegations

(GATE EE 2025)
  \end{itemize}




59. The cost function for a product in a firm is given by $5q^2$, where $q$ is the amount of production.  
The firm can sell the product at a market price of 50 per unit.  
The number of units to be produced by the firm such that the profit is maximized is

\begin{itemize}
\item[(A)] 5

\item[(B)] 10

\item[(C)] 15

\item[(D)] 25

(GATE EE 2025)
  \end{itemize}




60. Choose the most appropriate alternative from the options given below to complete the following sentence:  
\textit{Suresh’s dog is the one \rule{1cm}{0.15mm} was hurt in the stampede.}
\begin{itemize}
\item[(A)] that

\item[(B)] which

\item[(C)] who

\item[(D)] whom

(GATE EE 2025)
  \end{itemize}




% --------------------------
% GENERAL APTITUDE Q61–Q65
% --------------------------

61. Which of the following assertions are CORRECT?  
P: Adding 7 to each entry in a list adds 7 to the mean of the list.  
Q: Adding 7 to each entry in a list adds 7 to the standard deviation of the list.  
R: Doubling each entry in a list doubles the mean of the list.  
S: Doubling each entry in a list leaves the standard deviation of the list unchanged.

\begin{itemize}
\item[(A)] P, Q

\item[(B)] Q, R

\item[(C)] P, R

\item[(D)] R, S

(GATE EE 2025)
  \end{itemize}




62. An automobile plant contracted to buy shock absorbers from two suppliers X and Y.  
X supplies 60% and Y supplies 40% of the shock absorbers. All shock absorbers are subjected to a quality test.  
Of X’s shock absorbers, 96% are reliable. Of Y’s shock absorbers, 72% are reliable.  
The probability that a randomly chosen shock absorber, which is found to be reliable, is made by Y is

\begin{itemize}
\item[(A)] 0.288

\item[(B)] 0.334

\item[(C)] 0.667

\item[(D)] 0.720

(GATE EE 2025)
  \end{itemize}




63. A political party orders an arch for the entrance to the ground in which the annual convention is being held.  
The profile of the arch follows the equation $y = 2x - 0.1x^2$ where $y$ is the height of the arch in meters.  
The maximum possible height of the arch is

\begin{itemize}
\item[(A)] 8 meters

\item[(B)] 10 meters

\item[(C)] 12 meters

\item[(D)] 14 meters

(GATE EE 2025)
  \end{itemize}




64. \textit{Wanted Temporary, Part-time persons for the post of Field Interviewer to conduct personal interviews to collect and collate economic data. Requirements: High School-pass, must be available for Day, Evening and Saturday work. Transportation paid, expenses reimbursed.}


Which one of the following is the best inference from the above advertisement?
\begin{itemize}
\item[(A)] Gender-discriminatory

\item[(B)] Xenophobic

\item[(C)] Not designed to make the post attractive

\item[(D)] Not gender-discriminatory

(GATE EE 2025)
  \end{itemize}




65. Given the sequence of terms, AD, CG, FK, JP, the next term is

\begin{itemize}
\item[(A)] OV

\item[(B)] OW

\item[(C)] PV

\item[(D)] PW

(GATE EE 2025)
  \end{itemize}




\twocolumn[
\begin{center}
\textbf{END OF THE QUESTION PAPER}
\end{center}
]

% --------------------------
% Continue in same pattern...
% --------------------------

\end{document}


